\subsection{Sandbox und Training der KI Modelle}

Die beiden Data Science Use Cases, 
sowohl die Betrugserkennung bei KFZ-Schadensfällen
als auch die Prognose von Kündigern von privaten Krankenversicherungen,
sollen jeweils mit einem KI-Modell umgesetzt werden, 
welches speziell für diesen Anwendungsfall trainiert wird
\citep[Folien~11~\&~17]{nordakademieSzenarioVersicherung}.
Für die Datenarchitektur ergeben sich daraus
vier Analyseanforderungen: 
je Anwendungsfall einerseits das Training des KI-Modells
auf historischen Daten 
und andererseits der Einsatz des KI-Modells für 
aktuelle Datensätze.
Training und Einsatz der KI-Modelle bauen jeweils auf 
unterschiedlichen Architekturkomponenten auf.
Im folgenden wird zunächst der Teil des Trainings betrachtet.

Bei dem Training der KI-Modelle für die Beiden Use Cases
handelt es sich um unstrukturierte Datenanalysen.
Die Muster, die der Machine Learning Algorithmus
in den Daten findet, sind vorher noch nicht bekannt
und die Komplexität geht weit über standardmäßiges BI Reporting hinaus.
Trotzdem wird jeweils genau eine spezifische Analysefragestellung untersucht.

Für das Training soll als ergänzende Architekturkomponente 
eine Sandbox genutzt werden.
In der Sandbox werden erstens lediglich Kopien der Daten abgelegt.
Und zweitens haben die Analyseexperten und -expertinnen  
vollen Zugriff auf die Daten und können diese manipulieren,
ohne damit den Datenbestand im Data Warehouse oder Data Lake zu verändern.
Diese Eigenschaften sind hier von Vorteil
um bei der explorativen und unstrukturierten Datenanalyse
uneingeschränkt vorgehen zu können.
Für die Use Cases können die Daten in der Sandbox intensiv 
und mehrmals in unterschiedlichen Arten und Weisen aufbereitet werden.
Und es können verschiedene ML Algorithmen ausprobiert werden,
bis das KI-Modell die gewünschte Qualität erreicht.
Zwischen den Iterationen kann die Sandbox immer wieder neu
aufgebaut werden, ohne, dass Daten verloren gehen.

Die Trennung der Sandbox von Data Warehouse und Data Lake 
sorgt zusätzlich dafür, dass letztere nicht durch das
Modell-Training belastet werden, was möglicherweise 
Performance-intensiv ist.
Je Anwendungsfall wird außerdem eine eigene Sandbox verwendet,
um die Datenbasen für die Projekte einzeln zusammenstellen und bearbeiten zu können.

Neben dem Data Warehouse und Data Lake werden für den ersten Use Case
(KFZ Betrugserkennung) Daten aus einer weiteren Datenquelle in die Sandbox geladen.
Dabei handelt es sich um eine externe Fraud Datenbank.
Dadurch wird sichergestellt, dass ausreichend viele Betrugsfälle
in den Trainingsdaten enthalten sind.
Da diese Daten allerdings nicht mit den Kunden der Allianz
in Verbindung stehen, sondern nur für das Modell-Training notwendig sind,
werden sie direkt in die Sandbox geladen.
Eine dauerhafte Speicherung im eigenen Data Warehouse wäre überflüssig.

Für das initiale Training der KI-Modelle, 
werden die Daten manuell und unregelmäßig von den Analyseexperten und -expertinnen
in die Sandboxen geladen.
Aufgrund von Phänomenen wie Data Drift und Concept Drift
ist es jedoch häufig unabdingbar die Modelle nach zu trainieren (sog. ''Retraining''),
damit sie nicht an Performance verlieren \cite{ghosh2022}.
Deshalb ist in Abb. \ref{fig:datenarchitektur-regional} an allen gerichteten Kanten,
die in die Sandbox führen,
alternativ die aktuelle Datenbereitstellung in Batchverarbeitung vermerkt.
Diese soll das Retraining in regelmäßigen Abständen repräsentieren. 

% TODO: Diagram 2 Sandboxen gefächert?
% Trainingsdaten aus Sandbox archivieren? (Archiv/Data Lake)